\chapter{The Subcanonicity Countermodel and the Representable Repair}\label{ch06:main}

Chapter~\ref{ch05:main} established that the legacy topology is exactly the
jointly-surjective observational topology and that the modalities intended
to ride on the ACBE and RPM clauses are trivial. One question about the
legacy site remains, and it is the question that determines how legacy
\emph{objects}---as opposed to legacy covers---must be handled by any
successor: is the legacy topology subcanonical? If every representable
presheaf $yX$ is a sheaf, then the objects of $\mathrm{Noetica}$ pass into
the sheaf topos unmodified, carrying their morphisms and their identities
with them, and every downstream artifact may treat ``the legacy object $X$''
and ``the sheaf $yX$'' as the same thing. If subcanonicity fails, or merely
fails to be provable, that identification is illegitimate, and a repair with
its own discipline is required.

This chapter shows that subcanonicity is \emph{not derivable} from the
registered axioms, by exhibiting a complete, finite, elementary
countermodel: the category of sets of size at most two with injections,
carrying its jointly-surjective coverage. In that model the representable of
a one-element set fails the sheaf condition against the two-point cover of a
two-element set, by a computation short enough to display in full and
instructive enough to reward the display: the matching-family set is a
singleton, the global-section set is empty, and the gluing map from the
empty set to the singleton has nowhere to send nothing. We then delimit the
exact logical status of the result---non-derivability, not universal
non-subcanonicity---and present the standard repair: the distinguished
legacy object is represented in the sheaf topos by the associated sheaf
$a_{\mathrm{legacy}}(yX)$, never by silently treating $yX$ as a sheaf
\cite{stacks00WO}. Because the countermodel is fully computable, we can go
further than the repair's statement and compute exactly what the repair does
in the model---the answer is drastic, and knowing it is drastic is itself
valuable calibration. The chapter closes with extended remarks on
countermodel methodology in axiomatic auditing, since the technique used
here---small finite model, exhaustive verification, sharp scope
statement---is a template the program relies on repeatedly.

\begin{boundary}\label{ch06:bnd:status}
Status discipline. The countermodel argument originates in the registered
mathematical red-team report of Project Coherence (registered under SHA-256
\texttt{03AEC81BF9F6AA12376C2DF78E6CC410EB636589DA37D2E2534460132BE720C6})
and is restated here with full proofs; its status, like that of every finite
result in this monograph, is \textsc{informal proof present}---elementary,
finite, and checkable, but not peer-reviewed and not formalized in a proof
assistant. The extended computations of
Section~\ref{ch06:sec:repaircomputed} (the explicit identification of the
associated sheaves in the countermodel) are tutorial-level verifications
presented in this monograph with informal proofs; they are not among the
registered certificate checks, and nothing promotes them by their appearance
here. The chapter makes no empirical, causal, safety, or product claim.
\end{boundary}

\section{What subcanonicity buys}\label{ch06:sec:buys}

A topology $J$ on a small category $\cC$ is \emph{subcanonical} when every
representable presheaf is a $J$-sheaf
(Definition~\ref{ch04:def:sheaf}; \cite{maclane,johnstone}). The definition
is one line; its consequences are the reason anyone cares, and it is worth
listing them with some precision, because each one is something the legacy
construction implicitly wanted and---as this chapter shows---was not
entitled to assume.

First, \emph{the site embeds in its topos}. If $J$ is subcanonical, the
Yoneda embedding factors through sheaves,
\[
y \;:\; \cC \longrightarrow \Sh(\cC,J) \subseteq \PSh(\cC),
\]
and remains full and faithful: for all objects $X,Y$,
\[
\Hom_{\Sh(\cC,J)}(yX,\,yY)\;=\;\Hom_{\PSh(\cC)}(yX,\,yY)\;\cong\;\Hom_{\cC}(X,Y),
\]
the first equality because $\Sh$ is a full subcategory, the second by the
Yoneda lemma. The objects of the site appear in the topos \emph{as
themselves}: same morphisms, no new morphisms, no identifications. For a
program whose objects carry registered identities, this is the property that
lets ``the object $X$'' name one thing on both sides of the site/topos
boundary.

Second, \emph{evaluation is representable}. For any sheaf $F$, the Yoneda
lemma gives $\Hom(yX,F)\cong F(X)$; under subcanonicity this computes
sections of sheaves by maps in the topos from an object that is itself a
sheaf. Diagrammatic reasoning about sections---gluing arguments, naturality
arguments, limit computations---can then be conducted entirely inside
$\Sh(\cC,J)$.

Third, \emph{covers become genuine colimit statements}. A covering sieve
$S\in J(X)$ with $yX$ a sheaf exhibits $yX$ as a colimit of the
representables over $S$ in the sheaf category; covering families become
effective-epimorphic families. The topology's covers acquire an exactness
meaning, not merely a bookkeeping one \cite{maclane,johnstone}.

Fourth, there is a largest subcanonical topology on any small category, the
\emph{canonical topology}, whose covers are the universally
effective-epimorphic sieves (Remark~\ref{ch04:rem:canonical}). Subcanonicity
of $J$ is precisely the containment of $J$ in the canonical topology.
Asserting subcanonicity is therefore asserting a comparison between the
given coverage and an intrinsic invariant of the category---and comparisons
require proofs. The legacy source supplied none, and the next section
explains why none was forthcoming.

\section{Why the legacy axioms do not derive subcanonicity}
\label{ch06:sec:whynot}

Unfold what subcanonicity demands of the jointly-surjective coverage. Let
$\{f_i:U_i\to X\}$ be a jointly surjective family and $Y$ any object. The
sheaf condition for $yY$ along (the sieve generated by) this family says:
for every family of morphisms $g_i:U_i\to Y$ compatible on all overlaps,
there is a unique morphism $g:X\to Y$ with $g\circ f_i=g_i$. That is: the
family is \emph{effective-epimorphic}, and since covers pull back, it must
be so universally.

Joint surjectivity is a statement about \emph{elements}: every point of $X$
is hit. Effective epimorphy is a statement about the \emph{morphism supply}
of the category: the function one would define piecewise---send
$x=f_i(u)$ to $g_i(u)$, well defined by compatibility---must actually be an
arrow of $\cC$. In the category of all finite sets and all functions, it
always is: the piecewise definition is a function, functions are arrows, and
jointly surjective families are universally effective-epimorphic; the
jointly-surjective coverage on that category is subcanonical. The registered
axioms, however, do not say that $\mathrm{Noetica}$ has all functions as
morphisms. They say its morphisms are generated by designated
admissibility-relation maps and closed under pullback---a \emph{restricted}
morphism supply, of unspecified extent. Nothing in the axioms requires the
piecewise-glued function to be admissible. The gap between
elementwise covering and morphism-level gluing is exactly the gap the
axioms leave open, and a derivation of subcanonicity would have to close it
for every admissible instantiation. The countermodel shows it cannot be
closed: there is an instantiation---finite concrete sets, a
composition-closed and pullback-closed class of structure maps, the
jointly-surjective coverage---in which a glued arrow fails to exist at the
very first nontrivial cover.

One more framing before the construction. The failure mode has a
recognizable signature, previewed in Example~\ref{ch04:ex:nonposet}:
\emph{vacuous compatibility plus missing gluing}. When cover pieces have
small or empty overlaps, the compatibility condition on a would-be glued
family degenerates---every family is compatible---while the gluing
requirement stays fully strong. Categories with restricted morphism supply
are precisely where this bites, because the compatibility condition is the
only brake on the demand side, and the morphism supply is the only source on
the offer side. The countermodel is engineered so that the brake is fully
released and the source is empty.

\section{The countermodel}\label{ch06:sec:countermodel}

\begin{definition}[The countermodel site]\label{ch06:def:inj2}
Fix a countable set $\mathcal{U}$ of atoms (so that the category is small)
and let $\mathcal{I}_{\le 2}$ be the category whose objects are the subsets
of $\mathcal{U}$ of size at most two, and whose morphisms are the
\emph{injections} between them, with ordinary composition of functions.
Declare a sieve $S$ on an object $X$ to be covering when the family of its
members is jointly surjective:
$\bigcup_{f\in S}\im(f)=X$. Write $J_{\mathrm{js}}$ for this assignment.
\end{definition}

Everything about $\mathcal{I}_{\le 2}$ is finite and enumerable at each
object, and it is useful to have the counts explicit. For objects $D$, $X$
with $|D|=d$ and $|X|=n$ (each in $\{0,1,2\}$), the morphism counts are:
\[
|\Hom(D,X)| \;=\;
\begin{cases}
1 & d=0,\\
n & d=1,\\
n(n-1) & d=2,
\end{cases}
\]
so in particular $\Hom(D,X)=\varnothing$ whenever $d>n$: there is no
injection from a two-element set to a one-element set (pigeonhole), and this
single emptiness is the engine of the entire countermodel. Composites of
injections are injections; identities are injections; and pullbacks exist
and are computed as in $\Set$: the pullback of two injections
$U\to X\leftarrow V$ is (isomorphic to) the intersection of their images,
which has size at most two, so $\mathcal{I}_{\le 2}$ is closed under the
pullbacks the site machinery uses. Up to isomorphism the category has three
objects---$\varnothing$, a singleton, a doubleton---and readers who prefer a
skeleton may work with $0=\varnothing$, $1=\{1\}$, $2=\{1,2\}$ throughout;
we keep named sets $\{a\}$, $\{b\}$, $\{a,b\}$, $\{*\}$ for readability.

\begin{lemma}[Covering sieves contain all point inclusions]
\label{ch06:lem:points}
Let $S$ be a sieve on $X$ in $\mathcal{I}_{\le 2}$. Then $S$ is jointly
surjective if and only if for every $x\in X$ the point inclusion
$p_x:\{x\}\hookrightarrow X$ belongs to $S$. (For $X=\varnothing$ every
sieve, including the empty sieve, is jointly surjective.)
\end{lemma}

\begin{proof}
($\Leftarrow$) The images of the point inclusions are the singletons
$\{x\}$, whose union is $X$. ($\Rightarrow$) Let $x\in X$ and choose
$f:E\to X$ in $S$ with $x\in\im(f)$, say $f(e)=x$. Let $k:\{x\}\to E$ be the
injection sending $x$ to $e$. Then $f\circ k:\{x\}\to X$ is a member of $S$
by precomposition closure, and it is an injection from $\{x\}$ sending $x$
to $x$---but any two injections $\{x\}\to X$ with the same image point are
equal as functions, and $f\circ k$ has image $\{x\}$, so $f\circ k=p_x$.
Hence $p_x\in S$.
\end{proof}

\begin{proposition}[$J_{\mathrm{js}}$ is a Grothendieck topology]
\label{ch06:prop:topology}
The assignment $J_{\mathrm{js}}$ of Definition~\ref{ch06:def:inj2}
satisfies (J1), (J2), and (J3); hence
$(\mathcal{I}_{\le 2},J_{\mathrm{js}})$ is a site.
\end{proposition}

\begin{proof}
(J1) The maximal sieve on $X$ contains $\id_X$, whose image is $X$.

(J2) Let $S$ be covering on $X$ and $g:D\to X$ any arrow; we must show
$g^{*}S=\{h \mid g\circ h\in S\}$ is covering on $D$. Let $d\in D$. Then
$g(d)\in X$, so by Lemma~\ref{ch06:lem:points} the point inclusion
$p_{g(d)}:\{g(d)\}\to X$ lies in $S$. Consider the point inclusion
$p_d:\{d\}\to D$; the composite $g\circ p_d:\{d\}\to X$ is an injection with
image $\{g(d)\}$. As in the proof of Lemma~\ref{ch06:lem:points},
$g\circ p_d$ equals $p_{g(d)}$ precomposed with the bijection
$\{d\}\to\{g(d)\}$, hence lies in $S$ by closure. So $p_d\in g^{*}S$, for
every $d\in D$; by Lemma~\ref{ch06:lem:points}, $g^{*}S$ is covering. (If
$D=\varnothing$ there is nothing to check.)

(J3) Let $S$ be covering on $X$ and let $T$ be a sieve on $X$ such that
$f^{*}T$ is covering on $D$ for every $f:D\to X$ in $S$. Let $x\in X$; we
find a member of $T$ hitting $x$. By Lemma~\ref{ch06:lem:points},
$p_x\in S$; by hypothesis $p_x^{*}T$ is covering on $\{x\}$, so it contains
the point inclusion $\id_{\{x\}}$ (Lemma~\ref{ch06:lem:points} again, at the
object $\{x\}$). Membership $\id_{\{x\}}\in p_x^{*}T$ means
$p_x\circ\id=p_x\in T$, and $\im(p_x)\ni x$. Hence $T$ is jointly
surjective, i.e.\ covering.
\end{proof}

\begin{remark}[The countermodel instantiates the registered
axioms]\label{ch06:rem:instantiation}
To serve as a countermodel to \emph{derivability}, the site must satisfy
what the registered axioms demand, at the level at which a derivation would
have to operate (Boundary~\ref{ch05:bnd:schema}). The checklist: the objects
are finite concrete sets; the category is small; the morphisms form a
designated class of structure-preserving maps (injections) containing all
identities, closed under composition and under the pullbacks used by the
site machinery; the observational clause (O) is instantiated by the
jointly-surjective coverage just verified. Clauses (A) and (R) are
instantiated degenerately but legitimately: take the ACBE designation to be
the reflexive one (the identity at each context), and take the RPM data to
be $R(X)=X$ with $\eta_X=\varepsilon_X=\id_X$, a valid section-retraction
pair. By Theorem~\ref{ch05:thm:collapse} these clauses contribute only
maximal sieves in \emph{every} model, so the total legacy topology of this
model is exactly the $J_{\mathrm{js}}$ of
Proposition~\ref{ch06:prop:topology}. Where the registry's fuller
instantiation data (the forty-index object class, the specific admissibility
generators) are not reproduced here, that is because a derivation of
subcanonicity from the axioms could not lean on them either---no registered
axiom ties the index structure to the existence of glued morphisms, and a
proof from the axioms must survive every instantiation the axioms permit.
That is precisely what the next section refutes.
\end{remark}

\section{The failed gluing, in full}\label{ch06:sec:failure}

We now carry out the central computation with nothing suppressed. Fix two
atoms $a\ne b$ and a third atom $*$, giving objects $\{a\}$, $\{b\}$,
$\{a,b\}$, $\{*\}$ of $\mathcal{I}_{\le 2}$.

\begin{lemma}[The cover and its sieve]\label{ch06:lem:sieve}
Let $S_2$ be the sieve on $\{a,b\}$ generated by the two point inclusions
$p_a:\{a\}\hookrightarrow\{a,b\}$ and $p_b:\{b\}\hookrightarrow\{a,b\}$.
Then:
\begin{enumerate}
\item $S_2$ is covering for $J_{\mathrm{js}}$;
\item $S_2$ consists exactly of the arrows into $\{a,b\}$ whose domain has
size at most one---that is, all four injections from singletons (two with
image $\{a\}$, two with image $\{b\}$, allowing arbitrary singleton domains
from $\mathcal{U}$) at each singleton domain, together with the unique
arrow from $\varnothing$;
\item $S_2$ is not maximal: neither $\id_{\{a,b\}}$ nor the swap
automorphism belongs to it;
\item $S_2$ is the \emph{minimum} covering sieve on $\{a,b\}$: it is
contained in every covering sieve.
\end{enumerate}
\end{lemma}

\begin{proof}
(1) The images $\{a\}$ and $\{b\}$ of the generators union to $\{a,b\}$.

(2) Membership in the generated sieve means factoring through $p_a$ or
$p_b$. Any arrow $f:D\to\{a,b\}$ with $|D|\le 1$ has image contained in
$\{a\}$ or in $\{b\}$ (its image has at most one point), hence factors
through the corresponding point inclusion: for $D=\varnothing$ via the empty
injection into either singleton, and for $D=\{d\}$ via the bijection
$\{d\}\to\{f(d)\}$ followed by $p_{f(d)}$. Conversely any arrow factoring
through a singleton has domain of size at most one, injections being
injective. So the members are exactly the arrows from objects of size at
most one.

(3) The two automorphisms of $\{a,b\}$ have two-element domain, hence are
not members by (2). (Equivalently: were either present,
Lemma~\ref{ch05:lem:iso} would make $S_2$ maximal, and (2) shows it is
not.)

(4) By Lemma~\ref{ch06:lem:points} every covering sieve on $\{a,b\}$
contains $p_a$ and $p_b$, hence contains the sieve they generate.
\end{proof}

\begin{theorem}[The representable $y\{*\}$ is not a sheaf]
\label{ch06:thm:notsheaf}
Consider the representable presheaf $y\{*\}=\Hom_{\mathcal{I}_{\le 2}}(-,\{*\})$
and the covering sieve $S_2$ on $\{a,b\}$. Then:
\begin{enumerate}
\item the matching-family set $\mathrm{Match}(S_2,\,y\{*\})$ has exactly one
element;
\item the section set $y\{*\}(\{a,b\})$ is empty;
\item hence the gluing map
$y\{*\}(\{a,b\})\to\mathrm{Match}(S_2,\,y\{*\})$ is injective but not
surjective: $y\{*\}$ is separated for $S_2$ but has a matching family with
no amalgamation, and is therefore not a $J_{\mathrm{js}}$-sheaf.
\end{enumerate}
Consequently $(\mathcal{I}_{\le 2},J_{\mathrm{js}})$ is not subcanonical.
\end{theorem}

\begin{proof}
(1) A matching family assigns to each $f:D\to\{a,b\}$ in $S_2$ an element
$x_f\in y\{*\}(D)=\Hom(D,\{*\})$, naturally. By
Lemma~\ref{ch06:lem:sieve}(2), every member has $|D|\le 1$, and for such $D$
the set $\Hom(D,\{*\})$ is a \emph{singleton}: the empty injection when
$D=\varnothing$, the unique bijection $D\to\{*\}$ when $|D|=1$. So there is
exactly one candidate assignment---$x_f$ is the unique element of its
singleton---and no freedom anywhere. It remains to observe that the
candidate is in fact matching: for any $h:E\to D$, both $x_{f\circ h}$ and
$x_f|_h=x_f\circ h$ are elements of $\Hom(E,\{*\})$, which is again a
singleton (as $|E|\le|D|\le 1$), so they are equal. Hence
$\mathrm{Match}(S_2,y\{*\})$ is a singleton: the compatibility conditions
are all vacuously satisfied by forced data. This is the ``vacuous
compatibility'' half of the signature announced in
Section~\ref{ch06:sec:whynot}: the pieces $\{a\}$ and $\{b\}$ overlap only
over $\varnothing$, where $y\{*\}$ holds no information, so the cover
imposes no constraint that could rule the family out.

(2) $y\{*\}(\{a,b\})=\Hom(\{a,b\},\{*\})$ is the set of injections from a
two-element set to a one-element set, which is empty by pigeonhole. This is
the ``missing gluing'' half: the would-be amalgamation---the function
sending both $a$ and $b$ to $*$---exists perfectly well as a function, but
it is not injective, so it is not an arrow of $\mathcal{I}_{\le 2}$. The
category's morphism supply stops exactly one arrow short of what the cover
demands.

(3) The gluing map is a map $\varnothing\to\{\text{pt}\}$; it is vacuously
injective and not surjective. A sheaf requires bijectivity for every
covering sieve; $S_2$ is covering by Lemma~\ref{ch06:lem:sieve}(1); so
$y\{*\}$ is not a sheaf, and subcanonicity fails for the site.
\end{proof}

\begin{corollary}[Non-derivability]\label{ch06:cor:nonderivable}
Subcanonicity of the legacy site is not derivable from the registered
axioms: there exists a model of the axioms
(Remark~\ref{ch06:rem:instantiation}) whose jointly-surjective coverage is
not subcanonical (Theorem~\ref{ch06:thm:notsheaf}). Hence every artifact
that treated a bare legacy representable as a sheaf was relying on an
unproved assumption, and the v7.4 corrigendum's retraction of that usage
(Section~\ref{ch05:sec:corrigendum}) is mandatory, not precautionary.
\end{corollary}

\begin{remark}[How little was needed]\label{ch06:rem:minimal}
It is worth tallying the ingredients, because the economy of the
countermodel is itself methodologically instructive
(Section~\ref{ch06:sec:methodology}). The argument used: one cover with two
pieces whose overlap carries no data; one target object small enough that
pigeonhole empties the crucial hom-set; and the injectivity constraint on
morphisms, which is what prevents the set-level gluing from being an arrow.
No infinite structure, no exotic category, no interaction with the ACBE/RPM
clauses (which contribute only maximal sieves and hence cannot rescue
gluing: a maximal sieve's gluing condition is automatic and so adds no
constraint that would exclude the offending matching family---nor any new
way to satisfy it). Three objects up to isomorphism suffice. When an
axiom system's intended theorem fails, it usually fails small, and the
audit practice of searching small models first
(Protocol~\ref{ch06:prot:audit}) is calibrated to exactly this experience.
\end{remark}

\section{The exact logical status}\label{ch06:sec:status}

Precision about what has and has not been proved is the difference between
an audit finding and a rhetorical flourish. Three statements must be held
apart.

\emph{First: what is proved.} The registered axioms do not entail
subcanonicity. This is a non-derivability result, established by the
standard model-theoretic route: a structure satisfying the axioms in which
the candidate theorem is false. Any purported proof of subcanonicity from
the registered axioms alone is unsound---it would prove a false statement
about $(\mathcal{I}_{\le 2},J_{\mathrm{js}})$.

\emph{Second: what is not proved.} Universal non-subcanonicity is not
established. Nothing here shows that the intended, fully instantiated
$\mathrm{Noetica}$---with its registered forty-index object class and its
actual admissibility generators---fails subcanonicity. That question is
simply open, and it is a question about a particular site, answerable only
by working with that site's registered data. The countermodel neither
answers it nor prejudges it.

\emph{Third: the independence observation.} The same axiom schema also
admits models in which subcanonicity \emph{holds}, so the statement is
independent of the schema: neither it nor its negation is derivable. A
witness is already in hand from Chapter~\ref{ch04:main}: take the site to
be a finite powerset poset with inclusion maps---inclusions are injections
between finite sets, identities are present, composition and pullback
(intersection) closure hold---with the jointly-surjective coverage, which
for inclusions is exactly the union coverage. There, every representable is
a sheaf: if every piece of a covering family lands in $W$, their union,
which is the covered object, lands in $W$; this is
Example~\ref{ch04:ex:representables} generalized, and it is the same
argument that proves the successor's context-factor topologies subcanonical
later in the monograph. The two models bracket the axioms from both sides,
and the bracket is the correct summary: \emph{the registered axioms
underdetermine subcanonicity}. Deriving it for an exact revised site is a
registered open obligation of the program, with a proof required; until
that obligation closes, the repair of the next section is the only
licensed handling of legacy representables.

\begin{finding}\label{ch06:fnd:status}
Subcanonicity is independent of the registered legacy axiom schema:
$(\mathcal{I}_{\le 2},J_{\mathrm{js}})$ satisfies the axioms and refutes
it; finite powerset posets with union coverage satisfy the axioms and
verify it. All prior text asserting or silently using legacy subcanonicity
is retracted to the conditional form: \emph{if} a revised exact site is
proved subcanonical, \emph{then} representables may be used as sheaves
there; otherwise the associated-sheaf repair applies.
\end{finding}

\section{The standard repair: $a_{\mathrm{legacy}}(yX)$}
\label{ch06:sec:repair}

The repair is classical and requires no invention; what it requires is
discipline in use. Let $a_{\mathrm{legacy}}$ denote sheafification for the
legacy site, and let $X$ be a distinguished legacy object. The
representable $yX$ is a presheaf in any case; its associated sheaf
\[
a_{\mathrm{legacy}}(yX)\;\in\;\Sh(\cC,\,J_{\mathrm{legacy}}),
\qquad
\eta_{yX}\;:\;yX\longrightarrow a_{\mathrm{legacy}}(yX),
\]
is the canonical---and, absent a subcanonicity proof, the only
licensed---avatar of $X$ in the legacy sheaf topos \cite{stacks00WO}. Three
properties govern its use.

\begin{proposition}[What the repair preserves]\label{ch06:prop:preserves}
For every sheaf $G\in\Sh(\cC,J_{\mathrm{legacy}})$, composition with the
unit induces a bijection
\[
\Hom_{\Sh}\bigl(a_{\mathrm{legacy}}(yX),\,G\bigr)
\;\cong\;
\Hom_{\PSh}(yX,\,G)
\;\cong\;
G(X),
\]
naturally in $G$. That is: mapping \emph{out of} the repaired object into
sheaves computes exactly evaluation at $X$, just as mapping out of $yX$
does.
\end{proposition}

\begin{proof}
The first bijection is the universal property of sheafification
(Definition~\ref{ch04:def:assoc}); the second is the Yoneda lemma. Both are
natural in $G$.
\end{proof}

So a pipeline that uses legacy objects only as \emph{probes}---evaluating
sheaves at $X$, forming matching data over covers of $X$, indexing
assessments by $X$---loses nothing under the repair: the repaired object
represents evaluation at $X$ on the entire sheaf category. What the repair
does \emph{not} preserve is the internal structure of the object itself:
its sections over other objects, its automorphisms in the topos, its
distinctness from other repaired objects. The countermodel makes this
concrete, and the computation is worth doing in full because it
demonstrates that the repair can be as far from innocuous as it is from
optional.

\subsection{The repair computed in the countermodel}
\label{ch06:sec:repaircomputed}

\begin{proposition}[Associated sheaves of representables in
$\mathcal{I}_{\le 2}$]\label{ch06:prop:computed}
In $(\mathcal{I}_{\le 2},J_{\mathrm{js}})$, write $\Set(-,X)$ for the
presheaf $D\mapsto\Set(D,X)$ of \emph{all} functions into the underlying
set of $X$, with restriction by precomposition. Then for every object $X$:
\begin{enumerate}
\item $\Set(-,X)$ is a $J_{\mathrm{js}}$-sheaf;
\item the evident inclusion $yX\hookrightarrow\Set(-,X)$ (injections are
functions) exhibits $\Set(-,X)$ as the associated sheaf of $yX$, i.e.\
$a(yX)\cong\Set(-,X)$ with the inclusion as unit;
\item in particular $a(y\{*\})\cong\mathbf{1}$, the terminal sheaf: the
repaired singleton representable has exactly one section over every
object.
\end{enumerate}
\end{proposition}

\begin{proof}
(1) We check the sheaf condition using the minimum covering sieves, which
exist at every object; by Lemma~\ref{ch06:lem:points} every covering sieve
on $D$ contains the sieve generated by the point inclusions of $D$, and a
presheaf gluing bijectively along a subsieve that is itself covering glues
bijectively along any larger covering sieve provided its verdicts are
consistent---here we simply check all cases directly. At $D=\varnothing$:
the empty sieve covers, and $\Set(\varnothing,X)$ is a singleton (the empty
function), as required (Lemma~\ref{ch04:lem:null}). At $|D|=1$: by
Lemma~\ref{ch06:lem:points} a covering sieve contains
$\id_D$ (the point inclusion of the unique point is the identity up to the
bijection of singletons; precomposition closure then yields $\id_D$ itself),
hence is maximal, and maximal sieves glue freely
(Example~\ref{ch04:ex:maximal}). At $|D|=2$, say $D=\{x,y\}$: a matching
family for the minimum sieve $S_{\min}(D)$ assigns to each singleton-domain
member an element of a $\Set(E,X)$ with $|E|=1$, naturally; unwinding as in
Theorem~\ref{ch06:thm:notsheaf}(1), the family is freely determined by its
values at the two point inclusions $p_x,p_y$, i.e.\ by a pair of functions
$\{x\}\to X$, $\{y\}\to X$, i.e.\ by an arbitrary function $D\to X$; and
the gluing map sends $g\in\Set(D,X)$ to exactly that pair. It is therefore
a bijection. Larger covering sieves on $D$ contain additionally arrows with
doubleton domain (isomorphisms onto $D$ or injections into $D$ from other
doubletons---the latter do not exist for $|D|=2$ except isomorphisms);
a matching family's values there are forced by naturality from the point
values, and the same bijection results.

(2) By Proposition~\ref{ch04:prop:uniqueness} it suffices to show the
inclusion $\iota:yX\to\Set(-,X)$ has the universal property: every map
$\alpha:yX\to G$ with $G$ a sheaf factors uniquely through $\iota$. Both
existence and uniqueness follow from the plus-construction description, but
the direct route is cleaner here: the presheaf $\Set(-,X)$ is generated
over $yX$ by gluing, in the precise sense that every section
$g\in\Set(D,X)$ restricts along the minimum covering sieve of $D$ to
sections lying in the image of $\iota$ (over domains of size at most one,
every function is injective or empty, hence \emph{is} an injection---this
is the pivotal smallness fact). Given $\alpha:yX\to G$, define
$\bar\alpha_D(g)$ as the amalgamation in $G(D)$ of the matching family
$\bigl(\alpha_E(g\circ h)\bigr)_{(h:E\to D)\in S_{\min}(D)}$---well
defined since $G$ is a sheaf and the family matches by naturality of
$\alpha$; uniqueness of amalgamations makes $\bar\alpha$ natural and forces
it on all of $\Set(-,X)$, and $\bar\alpha\circ\iota=\alpha$ because for
injective $g$ the family already has the amalgamation $\alpha_D(g)$. Any
other factorization agrees with $\bar\alpha$ on restrictions to the
covering sieve, hence everywhere, by separatedness of $G$.

(3) $\Set(D,\{*\})$ is a singleton for every $D$.
\end{proof}

\begin{remark}[Reading the computation]\label{ch06:rem:reading}
Proposition~\ref{ch06:prop:computed} says that in the countermodel,
sheafification of representables \emph{manufactures exactly the missing
non-injective maps}: the topology demanded function-like gluing, the
category supplied only injections, and the associated sheaf completes the
hom-presheaf $\mathrm{Inj}(-,X)$ to the full function presheaf
$\Set(-,X)$. Three consequences deserve explicit statement, because each
is a live risk for any successor that handles legacy objects casually.
First, \emph{objects can collapse}: all repaired singleton representables
become terminal, hence pairwise isomorphic---$a(y\{a\})\cong
a(y\{b\})\cong\mathbf{1}$ even though $\{a\}$ and $\{b\}$ are distinct
objects of the site. The repair functor need not be injective on objects up
to isomorphism. Second, \emph{morphisms can appear}: computing with
Proposition~\ref{ch06:prop:preserves},
$\Hom_{\Sh}(a(yX),a(yY))\cong a(yY)(X)=\Set(X,Y)$, which strictly contains
$\Hom_{\mathcal{I}_{\le 2}}(X,Y)$ whenever a non-injective function
$X\to Y$ exists. The repaired image of the site sees maps the site does not
have; the composite $X\mapsto a(yX)$ is faithful here but not full. Third,
\emph{units need not be surjective or bijective}: the unit
$yX\to\Set(-,X)$ is a proper inclusion at every doubleton evaluated at
doubletons. None of this is pathology; it is what sheafification is
\emph{for}---enforcing descent that the site's raw hom-structure does not
satisfy. But it means that under a non-subcanonical (or
not-known-subcanonical) topology, ``the object $X$'' and ``the sheaf
avatar of $X$'' are different mathematical citizens, and no identification
between them may be made silently.
\end{remark}

\subsection{Handling discipline for legacy objects}
\label{ch06:sec:discipline}

\begin{ruleenv}[Representable repair rule]\label{ch06:rule:repair}
Until an exact revised legacy site is proved subcanonical, every successor
artifact that places a distinguished legacy object $X$ inside the legacy
sheaf topos must use $a_{\mathrm{legacy}}(yX)$ and must record that it has
done so. Prohibited: silently treating $yX$ as a sheaf; transferring
site-level identities of objects to topos-level isomorphism classes
(Remark~\ref{ch06:rem:reading} shows repaired objects may become
isomorphic); inferring site-level morphism counts from topos-level
hom-sets (the repair may add morphisms). Permitted without further proof:
using $a_{\mathrm{legacy}}(yX)$ as the representing object for evaluation
at $X$ on sheaves (Proposition~\ref{ch06:prop:preserves}); forming its
covers, subobjects, and diagrams inside the sheaf topos.
\end{ruleenv}

\begin{invariant}[Identity lives at the registry, not in the topos]
\label{ch06:inv:identity}
The identity of a legacy object---its registered identifier and canonical
payload hash---is site-level and serialization-level data. It must travel
with the marked records of the successor architecture, never be
reconstructed from topos-level isomorphism class, and never be presumed
stable under sheafification. A categorical isomorphism is not an identity
receipt. (This invariant restates, for objects, the division of labor that
Chapter~\ref{ch05:main} established for covers: the topos computes
semantics; the registry carries identity and provenance.)
\end{invariant}

\begin{remark}[If subcanonicity is later proved]\label{ch06:rem:dissolve}
The repair is self-liquidating. If a future exact revised site is proved
subcanonical, then for that site every $yX$ is already a sheaf, the unit
$\eta_{yX}:yX\to a(yX)$ is an isomorphism
(Proposition~\ref{ch04:prop:localssheaves}), and every use of
$a_{\mathrm{legacy}}(yX)$ made under Rule~\ref{ch06:rule:repair} becomes
canonically isomorphic to the bare representable---no migration, no debt,
no stranded artifacts. This is the reason the rule mandates the repaired
form rather than waiting: the repaired form is correct in both futures,
the bare form in only one. Successors should prefer constructions that
remain valid when open obligations close in either direction; the repair
is the type specimen of that design principle.
\end{remark}

\section{Countermodel methodology in axiomatic auditing}
\label{ch06:sec:methodology}

The countermodel of this chapter is a two-page argument, but it carried
real weight in the program: it converted an implicit assumption into a
named open obligation, forced a corrigendum, and imposed a handling rule on
every downstream artifact. This closing section abstracts the technique,
because the program's audit posture treats countermodel construction as a
first-class tool, on par with proof and executable check, and the abstract
lessons transfer beyond this instance.

\emph{Countermodels refute derivability, which is the auditable claim.}
An assertion inside a formal program rarely fails by being false in the
intended model, which is usually rich and forgiving; it fails by being
underivable from what was actually written down---and underivability is
what a countermodel certifies. The audit question is never ``do we believe
subcanonicity?'' but ``does the registered text entail it?''. Model both
answers: exhibit a structure satisfying exactly the registered text where
the claim fails (refuting entailment), and where possible a second
structure where it holds (blocking the overcorrection to
``the axioms refute it''). The bracket of Finding~\ref{ch06:fnd:status} is
the complete deliverable; either half alone invites misuse.

\emph{Search small, and search where the demands are strongest and the
supplies weakest.} The heuristic that found this countermodel generalizes:
identify what the candidate theorem demands (here: glued morphisms over
covers) and what the axioms supply (here: an unspecified morphism class),
then construct the leanest structure that honors every axiom while
starving the demand. Finiteness is not a concession but a weapon---finite
models are exhaustively checkable, communicable in full, and immune to the
suspicion that the failure hides an error in an infinite limit. The size
cap of two in $\mathcal{I}_{\le 2}$ is exactly the point at which
pigeonhole first bites; no smaller cap yields a countermodel, which is
itself worth verifying as an exercise (Exercise~\ref{ch06:exr:minimality}).

\emph{A countermodel is a permanent regression test.} Because the model is
finite and concrete, its failing computation can be executed, not merely
believed: enumerate the sieve (Lemma~\ref{ch06:lem:sieve}), count the
matching families (one), count the sections (zero), compare. The program's
practice of registering such computations alongside certificates---as it
does for the collapse mechanisms of Chapter~\ref{ch05:main}---turns each
audit finding into a machine-checkable invariant that future revisions must
continue to pass or consciously retire. An axiom change that silently
restores the broken derivation will trip the executable form of the
countermodel before it reaches prose.

\emph{Countermodels bound repairs.} The countermodel does not merely
destroy; it specifies. Any future subcanonicity proof for a revised site
must visibly use some feature that $(\mathcal{I}_{\le 2},J_{\mathrm{js}})$
lacks---a richer morphism supply, a restricted cover class, or additional
axioms---and referees of such a proof should demand to see where the
countermodel is excluded. Conversely the repair adopted in the meantime
(Rule~\ref{ch06:rule:repair}) is exactly the construction the countermodel
shows to be necessary and the universal property shows to be sufficient
for the mapping-out uses the program actually has. This
demand--supply--repair triangle is the shape of a completed audit item.

\begin{protocol}[Countermodel audit protocol]\label{ch06:prot:audit}
For a claim $\Sigma$ asserted or assumed within an axiomatized artifact:
(1) isolate the registered axioms $A$ actually written, distinguishing
them from intended-model folklore; (2) restate $\Sigma$ as a demand on
structures satisfying $A$; (3) attempt derivation sketches to locate the
load-bearing gap; (4) search finite structures satisfying $A$ that starve
the gap, smallest first; (5) on success, verify $A$ and refute $\Sigma$ in
the model with complete proofs, and where feasible register an executable
form; (6) determine the complementary side---does some model satisfy
$\Sigma$?---to classify $\Sigma$ as refuted, derivable, or independent;
(7) issue the scope statement (non-derivability versus universal failure),
the corrigendum for affected text, and the interim handling rule; (8)
record the obligations that would license removing the rule.
Steps (5)--(7) are the deliverable; a countermodel without a scope
statement and a handling rule is an incident report without a fix.
\end{protocol}

\begin{boundary}\label{ch06:bnd:closing}
Applying Protocol~\ref{ch06:prot:audit} to the present chapter: $\Sigma$ is
legacy subcanonicity; the model refuting it is
$(\mathcal{I}_{\le 2},J_{\mathrm{js}})$ with the degenerate but licensed
clause instantiation of Remark~\ref{ch06:rem:instantiation}; the model
satisfying it is the powerset-poset instantiation of
Section~\ref{ch06:sec:status}; the classification is \emph{independent};
the corrigendum item is registered in v7.4; the handling rule is
Rule~\ref{ch06:rule:repair}; and the discharging obligation---a
subcanonicity proof for an exact revised site---remains \textsc{open} in
the registered obligation ledger. No part of this chapter licenses any
statement beyond that classification.
\end{boundary}

\begin{exercise}\label{ch06:exr:minimality}
Show that the size cap two is minimal: in the category of sets of size at
most one with injections, every jointly-surjective sieve on every object
is maximal (use Lemma~\ref{ch06:lem:points} and
Lemma~\ref{ch05:lem:iso}), hence the coverage is subcanonical there.
Conclude that every countermodel to subcanonicity built from injections
requires an object with at least two elements, and identify exactly which
step of Theorem~\ref{ch06:thm:notsheaf} uses it.
\end{exercise}

\begin{exercise}\label{ch06:exr:othersheaves}
Determine all sheaves on $(\mathcal{I}_{\le 2},J_{\mathrm{js}})$ up to
isomorphism in terms of their values at the skeleton objects $0,1,2$ and
the induced structure maps: show a sheaf is determined by a set
$F_1=F(1)$ together with an action of the two-element automorphism group
of $2$ on $F(2)\cong F_1\times F_1$ compatible with the projections, and
verify your description against Proposition~\ref{ch06:prop:computed}.
\end{exercise}

\begin{exercise}\label{ch06:exr:emptyrep}
Verify that $y\varnothing$ \emph{is} a sheaf on
$(\mathcal{I}_{\le 2},J_{\mathrm{js}})$: its value is a singleton at
$\varnothing$ and empty elsewhere; check the null cover, the singleton
covers, and the doubleton covers separately, and reconcile the doubleton
case with Theorem~\ref{ch06:thm:notsheaf} (why does the empty
matching-family set rescue $y\varnothing$ but not $y\{*\}$?).
\end{exercise}

\begin{exercise}\label{ch06:exr:functor}
Using Proposition~\ref{ch06:prop:preserves}, show that in any site the
assignment $X\mapsto a(yX)$ extends to a functor $\cC\to\Sh(\cC,J)$, and
that it is full and faithful if and only if $J$ is subcanonical. Then
locate the failure of fullness in the countermodel explicitly: exhibit the
map $a(y\{a,b\})\to a(y\{a,b\})$ corresponding under
Remark~\ref{ch06:rem:reading} to the constant function and show it is not
in the image of $\Hom_{\mathcal{I}_{\le 2}}(\{a,b\},\{a,b\})$.
\end{exercise}
